\documentclass{article}
\usepackage[utf8]{inputenc}
\usepackage[T1]{fontenc}
\usepackage[margin=2cm]{geometry}
\usepackage{hyperref}
\sloppy
\begin{document}
\section*{Historical source}
Encyclopaedia Britannica scan, source row 1466. Parse 3.5 extracted the text. Super 3.5 VL corrected it from an enlarged crop of each block. It can still contain recognition errors.
Page headers, footers, and page numbers are removed.

\subsection*{Untitled continuation (E0122)}
\textit{Continues an entry from an earlier page.}

Effaying weighed, and made of paper not impregnated with alum. It is to be dried at first with a gentle heat, but afterwards exposed for five minutes to a heat of 100°. On cooling, it is to be weighed together with the filter; the known weight of which must afterwards be subtracted. The sediment is both washed in a bottle for a filter when once impregnated with saline matter cannot be freed from it again without great difficulty, especially if an interval of some hours intervenes.

The alkali made use of in Mr Bergman's experiments, was that of sola saturated with aerial acid. His phlogisticated alkali is made by deflagrating equal weights of pure nitre and cream of tartar intimately mixed together; the residuum is the common white flux. Half an ounce of this is dissolved in half a quadrans of distilled water. To this he adds, in a digesting heat, two ounces of Prussian blue, carefully avoiding such an effervescence as may throw any thing over, which easily happens if the quantity be too large. The pigment soon loses its beautiful blue colour, growing not red but black; which shows that a decomposition has taken place. The Prussian blue used in his experiments contained in 100 parts only 23 of the pigment and 77 of the clay; so that if we employ the blue made without any alum, 221 grains of it will saturate the half ounce of alkaline salt more completely than the two ounces of the kind already described. But in whatever manner the operation is performed, after the addition of the last quantity, the whole must be exposed to a stronger digesting heat, and stirred with a wooden spatula. If the liquor be too much diminished by evaporation, the defect must be supplied by adding more water. When the liquor becomes clear, the residuum must be collected upon filtering paper, and gradually washed with warm water until all the soluble part is extracted; when, if the operation has been properly conducted, the filtered liquor amounts to a whole quadrans, of a brownish yellow colour, and so well saturated with colouring matter, that it does not change the colour of paper tinged with Brazil wood. This lixivium, however, contains a small quantity of Prussian blue, about 4 lb. to a cwt. of the alkali. These should be previously separated by an acid, or, which is better, corrected by subtracting from the weight of the sediment 16 ounces for each quadrans of the lixivium. When we wish to examine the colour of the precipitate exactly, however, the lixivium we employ must necessarily be well depurated; for by neglecting this precaution we may easily persuade ourselves that any metal precipitated by the lixivium has a blue colour. When we only wish to ascertain the weight, the lixivium, having the small proportion of Prussian blue intermixed, may be employed; but still the proper correction must ultimately be made use of; for the precipitating acid is wont to impair the qualities of the lixivium, and even to destroy them altogether, especially in a warm temperature. Calcareous earth, whether in its mild or caustic state, is also capable of abstracting a coloured substance from iron and other metals.

In the precipitation of metals by metals, it is to be observed, that the acid of the solution ought to be somewhat predominant; but any considerable excess must be corrected occasionally, either by alkali, water, or spirit of wine.

In the following experiments an essay cwt. was always employed, unless where it is expressly mentioned otherwise; conclusions sufficiently accurate may indeed be obtained from 25 lb. may sometimes from smaller quantities. In these cases our author mentions the usual quantity; applying to them those formulae of calculation which are founded on the mutual proportions of the proximate principles constituting metallic faults. By an easy fublitution, the same formulae may be used by those who employ ½ or ¼ cwt. We now come to describe the method of effaying the ores of the particular metals.

1. Ores of Gold. This metal occurs in the bowels of the earth native, poffelling a complete metallic form, although in general the small particles of it are so in-

terfected in various matrices, that they are entirely invisible. It is also found mineralized, or united with sulphur by means of iron or some other metal. These two species of ore we shall consider separately.

Native gold is very seldom, if ever, free from heterogeneous matters; the most usual mixtures are copper, silver, and sometimes iron. The first of these remains in the menstruum, and may be separately collected by diffusing the gold in aqua-regia, and precipitating it by martial vitriol: the second falls during the solution, yielding a salted silver; which, being washed and dried, shows the weight of the silver contained; and the iron may be discovered by phlogisticated alkali. The precipitate occasioned by martial vitriol is pure gold in its metallic state, but very subtilely divided, and therefore its weight requires no correction.

divided, and therefore its weight requires no correction. Hence it appears how small a portion of gold inherent in the ores of other metals may be extracted; besides, a solution containing the most minute particle of gold instantly produces the purple precipitate of Cassius, with a solution of tin properly prepared.

As to the ore which contains gold adhering to and surrounded by flony particles, r. We must reduce a determined weight to an impalpable powder by triture and elutriation. Then let the powder, weighed a second time, be boiled in aqua regia, as long as any thing is taken up by the menstruum; after which, let the exhausted ore, well washed, be collected, exsiccated to ignition, and weighed. Let the clear solution (the colour of which, in some degree, affords a method of judging) be precipitated in the usual way by martial vitriol; the precipitate well washed, dried, and weighed, shows the gold, which, added to the weight of the exhausted ore, ought to be equal to the original weight, unless somewhat has been dispersed by the pulverisation, or unless some of the matrice has entered the menstruum. The former of these is discovered by comparing the weights before and after pulverisation; the latter by precipitants.

When grains of gold are mixed with loose earthy particles, they are sometimes easily separated by mechanical application of water.

When the metal is mineralised by sulphur, as in the golden pyrites, let one or more essay ewts. reduced to powder be gently boiled in the nitrous acid, or rather sulphur, digested in a heat of 50°—80°, left the sulphur should be destroyed. It is even necessary to employ a more gentle heat for this purpose, that the sulphureous par-


\section*{Source record}
\begin{description}
\item[Scan ID] \texttt{britannica:309f765b1ecc83e0a7f7762e44b4a927a7555443:train:1466}
\item[Image file] \texttt{page-1466-bb41810e6962.jpg}
\item[Image SHA-256] \texttt{bb41810e69624a12ba0855ee98b3a60633979581d1be11f1b33bdaf7f1fba919}
\item[Dataset] NationalLibraryOfScotland/encyclopaedia\_britannica\_illustrated at revision \texttt{309f765b1ecc83e0a7f7762e44b4a927a7555443} (CC0-1.0)
\item[Extraction] model \texttt{nvidia/nemotron-3.5-parse}, raw reply \texttt{extract/raw/page-1466-bb41810e6962.json}
\item[Raw reply SHA-256] \texttt{8730e1ecd921bd501c1ff29253449764745e86bd558d5cc11b3e54f26b4820a0}
\item[Correction] model \texttt{nvidia/Nemotron-3.5-Super-VL}, block crops, raw replies \texttt{correct/raw/crop\_nothink/page-1466-bb41810e6962-b*.json}; super\_vl: 11
\item[Running header] E S S [ 729 ] E S S
\end{description}
\end{document}
